\chapter{Medical kits}

\section*{Definition and Overview}
A medical kit, commonly referred to as a first aid kit, is a structured collection of medical supplies, dressings, and medications designed to provide immediate, provisional treatment for acute injuries, illnesses, and environmental emergencies. The composition of a medical kit varies based on its intended environment (such as the home, workplace, automotive, international travel, or remote wilderness) and the training level of the user. A well-stocked, easily accessible, and regularly maintained medical kit is a cornerstone of emergency preparedness. It allows individuals to stabilise acute conditions, prevent minor wounds from becoming infected, and manage self-limiting symptoms before professional medical assistance can be obtained.

\section*{Epidemiology}
Injuries represent a major public health challenge and are a leading cause of emergency department presentations and hospitalisations worldwide. In Australia, millions of minor injuries occur annually in homes, schools, workplaces, and during recreational activities. Safe Work Australia reports that workplace injuries and illnesses require immediate on-site first aid, and having compliant first aid facilities and kits is a legislative requirement for all employers under the Work Health and Safety (WHS) Act. In travel medicine, studies indicate that up to 50\% of international travellers experience some health-related issue during their trip, highlighting the clinical value of carrying a tailored travel medical kit to manage conditions like travellers' diarrhoea, minor trauma, and insect-borne illnesses.

\section*{Aetiology and Pathophysiology}
The necessity of a medical kit arises from the immediate physiological demands of acute trauma, sudden illness, or environmental exposures that require rapid intervention to prevent clinical deterioration.
\begin{itemize}
    \item \textbf{Soft tissue trauma:} Abrasions, lacerations, punctures, and burns disrupt the skin barrier, risking haemorrhage, bacterial contamination, and subsequent localised or systemic infection.
    \item \textbf{Musculoskeletal injuries:} Sprains, strains, fractures, and dislocations cause pain, inflammation, and joint instability, requiring immobilization to prevent secondary neurovascular damage.
    \item \textbf{Allergic reactions and anaphylaxis:} Systemic IgE-mediated hypersensitivity reactions to insect venoms, foods, or medications can trigger rapid airway oedema and cardiovascular collapse (anaphylactic shock).
    \item \textbf{Environmental extremes:} Prolonged exposure to heat or cold can lead to hyperthermia, hypothermia, dehydration, or severe sunburn.
    \item \textbf{Minor systemic illnesses:} Gastrointestinal disturbances (such as diarrhoea and vomiting), minor infections, motion sickness, and mild-to-moderate pain.
\end{itemize}

\section*{Clinical Features}
Determining when and how to use the contents of a medical kit requires recognising the clinical signs of common minor injuries and systemic emergencies.
\begin{itemize}
    \item \textbf{Active haemorrhage:} Bleeding from a wound. Capillary bleeding is slow and oozing; venous bleeding is dark red and flows steadily; arterial bleeding is bright red and spurts, representing a life-threatening emergency.
    \item \textbf{Localised inflammation:} Redness, swelling, warmth, and pain at the site of a musculoskeletal injury or minor burn.
    \item \textbf{Dehydration:} Dry mouth, decreased urine output, dizziness, and headache, particularly following prolonged heat exposure or gastrointestinal illness.
    \item \textbf{Allergic signs:} Urticaria (hives), pruritus (itching), and localised swelling. Severe signs include stridor, wheezing, swelling of the tongue, and hypotension.
    \item \textbf{Wound infection:} Worsening pain, expanding redness, warmth, swelling, and purulent discharge (pus) developing 24 to 72 hours after an injury.
\end{itemize}

\section*{Differential Diagnosis}
A critical aspect of first aid is differentiating minor conditions that can be managed safely on-site using a medical kit from severe pathologies requiring urgent professional medical care.
\begin{itemize}
    \item \textbf{Minor laceration vs. deep structural injury:} A superficial cut can be cleansed and dressed, whereas deep wounds with tendon, nerve, or major vascular involvement require surgical repair.
    \item \textbf{Simple joint sprain vs. bone fracture:} A sprain involves ligamentous stretching, whereas a fracture involves structural bone disruption, often presenting with deformity, bony crepitus, and inability to bear weight, requiring X-ray evaluation.
    \item \textbf{Mild allergic reaction vs. anaphylaxis:} Mild reactions are limited to the skin or mucosal surfaces, whereas anaphylaxis is a life-threatening multi-system emergency involving respiratory or cardiovascular compromise.
    \item \textbf{Uncomplicated travellers' diarrhoea vs. systemic dysentery:} Uncomplicated diarrhoea is watery and self-limiting, whereas dysentery presents with high fever, severe abdominal pain, and bloody stools, requiring medical evaluation and antibiotics.
\end{itemize}

\section*{When to Seek Medical Care}
First aid is an initial, temporary measure. Knowing when to escalate care to emergency services or a general practitioner is vital.

\textbf{Call 000 (or local emergency services) for:}
\begin{itemize}
    \item Severe, uncontrollable bleeding that does not stop with direct pressure.
    \item Anaphylaxis, characterised by difficulty breathing, swelling of the tongue or throat, or collapse (administer an adrenaline autoinjector immediately if available).
    \item Severe chest pain, pressure, or tightness, or a suspected myocardial infarction.
    \item Head injury accompanied by loss of consciousness, persistent vomiting, confusion, or pupillary asymmetry.
    \item Large, deep, or chemical burns, or any burns involving the face, hands, feet, or groin.
    \item Fractures with visible bone deformity, open wounds over the fracture site, or loss of distal pulses.
\end{itemize}

\section*{Diagnosis and Investigations}
On-site assessment of a patient requiring first aid relies on a rapid clinical survey rather than diagnostic tests.
\begin{itemize}
    \item \textbf{DRSABCD protocol:} The standardised primary survey: Danger (ensure safety of self and others), Response (check for consciousness), Send for help (call 000), Airway (clear and open), Breathing (look, listen, feel), CPR (commence chest compressions and rescue breaths if not breathing), Defibrillation (attach an Automated External Defibrillator if available).
    \item \textbf{Secondary survey:} A systematic head-to-toe examination to identify bleeding, fractures, or medical signs once life-threatening conditions are stabilised.
    \item \textbf{Wound assessment:} Checking the wound for depth, contamination, presence of foreign bodies, and assessing distal circulation, sensation, and movement.
\end{itemize}

\section*{Management}
A medical kit should be systematically organised, stored in a durable, water-resistant container, and checked regularly.
\begin{itemize}
    \item \textbf{Wound cleansing and antiseptics:} Sterile saline tubes or packets for wound irrigation, antiseptic wipes or solution (e.g., chlorhexidine or povidone-iodine), and alcohol-based hand sanitiser for the provider.
    \item \textbf{Dressings and bandages:} Sterile non-adherent pads, adhesive bandages (band-aids in various shapes), conforming crepe bandages (to hold dressings or support sprains), triangular bandages (for slings or splinting), sterile eye pads, and hypoallergenic adhesive tape.
    \item \textbf{Basic medical tools:} Sharp shears (to cut clothing and bandages), fine tweezers (to remove splinters or ticks), disposable nitrile gloves (multiple pairs to prevent cross-contamination), and a digital thermometer.
    \item \textbf{Symptomatic medications:} Paracetamol and ibuprofen (for pain and fever), oral rehydration salts (to treat dehydration), antihistamines (for allergic reactions), and hydrocortisone cream (for insect bites or contact dermatitis).
    \item \textbf{Emergency rescue items:} An adrenaline autoinjector (EpiPen) if a family member has known anaphylaxis, a salbutamol inhaler and spacer for asthma, a thermal emergency space blanket (for hypothermia or shock), and a pocket resuscitation mask with a one-way valve.
    \item \textbf{Kit maintenance:} Perform an audit every six months to check medication expiry dates, replace used items, and ensure batteries in torches or devices are functional.
\end{itemize}

\section*{Complications}
Improper use of medical kit contents or delays in seeking professional care can lead to secondary complications.
\begin{itemize}
    \item \textbf{Wound infection:} Results from inadequate cleansing of contaminated wounds or using non-sterile dressings.
    \item \textbf{Ischaemic tissue damage:} Caused by applying crepe bandages, splints, or tourniquets too tightly, obstructing distal arterial blood flow.
    \item \textbf{Adverse drug events:} Using expired medications, administering incorrect dosages, or failing to identify patient drug allergies.
    \item \textbf{Delays in definitive care:} Attempting to self-manage severe, surgical, or progressive medical emergencies without consulting a doctor.
\end{itemize}

\section*{Prognosis and Outcomes}
The prognosis for minor injuries and illnesses managed promptly with a medical kit is excellent, with most resolving completely without scarring or long-term disability. Early wound cleansing and dressing reduces the incidence of secondary bacterial infections. In critical emergencies, such as massive haemorrhage or anaphylaxis, the immediate availability and correct use of tourniquets, pressure dressings, and adrenaline autoinjectors are directly associated with increased survival rates.

\section*{Prevention and Self-Care}
Preparation and education are key to preventing emergencies and managing them effectively.
\begin{itemize}
    \item \textbf{First aid training:} Complete and regularly renew a certified first aid and CPR course (such as Provide First Aid, HLTAID011) to build practical skills and confidence.
    \item \textbf{Tailor the kit:} Customise the medical kit's contents based on activities (e.g., add vinegar for marine stings if visiting coastal areas, or extra water purification tablets for remote trekking).
    \item \textbf{Store safely:} Keep the medical kit in a clearly marked, easily accessible location known to all family members or employees, but out of reach of young children.
    \item \textbf{Avoid extreme heat:} Store the kit below 25 degrees Celsius where possible, as excessive heat can degrade medications (such as adrenaline) and adhesive properties of dressings.
\end{itemize}

\section*{Sources and Further Reading}
The following reputable organisations provide guidelines on first aid, emergency management, and medical kit recommendations:
\begin{itemize}
    \item St John Ambulance Australia. \textit{First aid kits and training resources.} https://www.stjohn.org.au/
    \item Australian Red Cross. \textit{First aid guides and services.} https://www.redcross.org.au/
    \item Safe Work Australia. \textit{First aid in the workplace code of practice.} https://www.safeworkaustralia.gov.au/
    \item Healthdirect Australia. \textit{First aid kits.} https://www.healthdirect.gov.au/first-aid-kits
    \item Better Health Channel. \textit{First aid kit checklist.} https://www.betterhealth.vic.gov.au/health/conditionsandtreatments/first-aid-kits
\end{itemize}
